\begin{thema}{A}
\begin{erwthma}
    \item Έστω η συνάρτηση $ f(x)=x^{\nu},\ \nu\in\mathbb{N}-\{0,1\} $. Να αποδείξετε ότι η $ f $ είναι παραγωγίσιμη στο $ \mathbb{R} $ και ισχύει $ f'(x)=\nu x^{\nu-1} $, δηλαδή $ \left(x^\nu\right)'=\nu x^{\nu-1} $.
    \item Τι ονομάζουμε ρυθμό μεταβολής ενός ποσού $y=f(x)$ ως προς ένα ποσό $x$ σε ένα σημείο $x_0$?
    \item Πότε μια συνάρτηση $f$ ονομάζεται συνεχής σε ένα σημείο $x_0$ του πεδίου ορισμού της?
    \item Να χαρακτηρίσετε τις προτάσεις που ακολουθούν, γράφοντας στο τετράδιό σας, δίπλα στο γράμμα που αντιστοιχεί σε κάθε πρόταση, τη λέξη \textbf{Σωστό}, αν η πρόταση είναι σωστή, ή \textbf{Λάθος}, αν η πρόταση είναι λανθασμένη.
    \begin{alist}
        \item Αν $f$ συνεχής στο $[a, \beta]$ και παραγωγίσιμη στο $(a, \beta)$ με $f(a)=f(\beta)$, υπάρχει μια τουλάχιστον ρίζα της $f'$ στο $(a, \beta)$.
        \item Κάθε περιττή συνεχής συνάρτηση στο $\mathbb{R}$ έχει ρίζα το $x=0$.
        \item Αν μια συνάρτηση $f$ είναι περιττή και συνεχής, το εμβαδόν του χωρίου μεταξύ της $C_f$, του $x'x$ και των $x=-a, x=a$ ($a>0$) ισούται με $2\dint_0^a |f(x)| \d x$.
        \item Τα κρίσιμα σημεία εντοπίζονται και στα άκρα του πεδίου ορισμού κλειστού διαστήματος.
        \item Αν $\lim\limits_{x\to x_0} f(x) = \ell$, τότε η συνάρτηση ορίζεται υποχρεωτικά στο $x_0$.
    \end{alist}
    \item Αν $f''(x) > 0$ για κάθε $x \in \Delta$, η γραφική παράσταση της $f$ αναφορικά με την εφαπτομένη της:
    \begin{tasks}[label=\let\textdexiakeraia\relax\alph*.](2)
        \task Βρίσκεται χαμηλότερα, με εξαίρεση το σημείο επαφής.
        \task Βρίσκεται ψηλότερα, με εξαίρεση το σημείο επαφής.
        \task Διαπερνάται από αυτή σε κάθε σημείο.
        \task Συμπίπτει με αυτή.
    \end{tasks}
\end{erwthma}
\end{thema}
